\subsubsection*{3.3.4.}
回到 3.3.1 的较一般假设, 取 $n=2m+1$. 对 $X_s$ 的每个奇点 $x$, 仍以 $\pm\delta_x$ 表示 $R^nf_!\Lambda(m)$ 在 $S$ 上的截面, 它们是 $H^{n-1}(Y,\Lambda(m))$ 的本原商的生成元(记号同上)在复合态射
\begin{center}
\begin{tikzcd}[column sep=huge,row sep=huge]
H^{n-1}(Y,\Lambda(m))
  & H^{n-1}(\widetilde X_{s(x)},\Lambda(m))
      \arrow[l,"\sim"'] \arrow[r]
  & H^{n-1}(X^*_{s(x)},\Lambda(m)) \arrow[d] \\
H^0\bigl(S,R^nf_!\Lambda(m)\bigr)
  & H_c^n(X_s,\Lambda(m)) \arrow[l,"\sim"']
  & H^n_{\{x\}}(X_s,\Lambda(m)) \arrow[l]
\end{tikzcd}
\end{center}
下的像. 类 $\pm\delta_x$ 称为 $x$ 处的\underline{消灭圈}.

当 $S$ 为特质时, 类
\[
\pm\delta\in H_c^n(X_{\bar\eta},\Lambda(m))
\]
是 3.3.3 中的类
\[
\pm\delta\in
H^n_{\{x\}}\bigl(X_s,R\Psi_{\bar\eta}(\Lambda(m))\bigr).
\]
在 XIII, 2.1.8.4 下的像.

我们将同时证明以下两个结果.

\subsubsection*{命题 3.3.5.}
\underline{设 $S$ 为 Hensel 特质, $X/S$ 如 3.3.3. 则变差由下式给出:}
\[
\operatorname{Var}(\sigma)(a)
=(-1)^{m+1}\varepsilon_{b(x)}(\sigma)(a,\delta)\delta.
\]

\subsubsection*{命题 3.3.6.}
\underline{设 $S$ 严格局部且正则, $f:X\to S$ 为相对维数 $n=2m+1$ 的固有平坦态射, 使 $X_s$ 恰有一个非光滑点 $x$, 并假设它为非退化二次点. 还假设 $b(x)$ 为一个参数. 则 Galois 群在几何一般纤维的上同调上的作用为}
\[
\sigma(a)=a+(-1)^{m+1}
\varepsilon_{b(x)}(\sigma)(a,\delta_x)\delta_x.
\]